\section{검증 방법론}
\label{sec:methodology}

\subsection{재현 가능한 일곱 단계}

본 작업에서 ``EasyCrypt로 검증한다''는 말은 수학식을 별도 모델에서 다시
계산한다는 뜻이 아니다. pinned Jasmin 구현에서 추출된 실제 procedure와
HAETAE 명세의 수학 object 사이에 refinement 정리를 세우고, 그 정리들을 실제
호출 순서대로 합성한다는 뜻이다. 전체 절차는 다음 일곱 단계로 고정한다.

\begin{enumerate}
  \item \textbf{소스와 추출 경계 고정.}
        Jasmin source hash, mode~2 호출 상수와 대상 procedure 목록을 고정하고
        \proc{jasmin2ec}로 EasyCrypt procedure를 생성한다.
  \item \textbf{수학 명세 고정.}
        HAETAE 명세의 환·정수 식, decomposition, norm predicate와 byte/object
        경계를 실행 코드와 독립적으로 정의한다.
  \item \textbf{표현 관계 정의.}
        수학 polynomial·vector·signature object가 구현의 NTT/CRT 배열,
        고정폭 word와 byte slice에 어떻게 표현되는지 predicate로 명시한다.
  \item \textbf{실제 procedure의 국소 Hoare 정리.}
        추출된 procedure를 직접 호출하여 precondition에서 실제 postcondition을
        도출한다. 성공이나 원하는 출력 equality를 전제로 넣지 않는다.
  \item \textbf{제어흐름과 정확성 양식 분리.}
        loop invariant, failure/success branch, frame을 증명하고 partial
        correctness, losslessness, termination을 서로 다른 주장으로 관리한다.
  \item \textbf{호출 순서 합성.}
        KeyGen, Sign, Verify의 국소 정리를 실제 caller/callee binding, 동일
        transcript와 decoded-signature-object identity 아래 합성하여 제한된
        최종 승인 정리를 만든다.
  \item \textbf{기계적 검증과 증거 보존.}
        fresh compile, proof-hole·axiom scan, extraction drift, safety evidence와
        claim ledger를 검사하고 재현 가능한 로그를 남긴다.
\end{enumerate}

이 순서에서 앞 단계는 뒤 단계의 theorem precondition을 정당화한다. 앞 단계가
미완료이면 그 사실을 새 가정으로 숨기지 않고 최종 claim의 경계로 공개한다.

현재 7일 MINCORE 실행에서는 decoded signature object를 합성 경계로 고정한다.
HBZ/rANS의 기존 결과는 재사용 증거이지만, \(h\), metadata, padding, full parser를
새로 닫는 작업은 보류한다. 따라서 byte codec 방법은 후속 확장의 방법론이지
이번 headline theorem의 완료 조건이 아니다.

\subsection{소스 고정, 추출과 증명 경계}

검증 단위는 함수 이름이 비슷한 새 observation model이 아니라
\cref{sec:scope}에 열거한 pinned Jasmin procedure의 EasyCrypt 추출 결과이다.
각 최종 정리 또는 그 직접 하위 정리는 해당 generated procedure를
\texttt{proc} 본문에서 직접 호출해야 한다. proof-friendly wrapper를 사용할
경우에는 wrapper와 generated procedure 사이의 exact equivalence를 먼저
증명한다. 이 조건은 구현의 branch, word 연산 또는 memory update를 수학식으로
대체한 뒤 그 대체 모델만 증명하는 것을 막는다.

추출은 다음 두 경계를 구분한다.

\begin{itemize}
  \item \textbf{EasyCrypt 실행 경계:} generated procedure의 명령 의미에 대한
        Hoare 결과이다.
  \item \textbf{compiled Jasmin 경계:} 위 결과에 Jasmin compiler theorem을
        적용한 결과이며, 대상 procedure의 safety 증거를 추가로 요구한다.
\end{itemize}

따라서 safety가 확보되지 않은 정리는 ``extracted EasyCrypt model 수준''으로만
표기하고 assembly-level functional correctness로 올리지 않는다. compiler,
extraction, EasyCrypt kernel 및 representation bridge는
\cref{sec:artifact-plan}의 TCB에 기록한다.

\subsection{수학 명세와 구현 표현의 연결}

명세는 \(\Rq\)의 polynomial·vector와 무한정수 연산으로 기술하지만 구현은
고정 길이 배열과 word를 사용한다. 각 procedure precondition에는 적어도 다음
세 층을 동시에 둔다.

\[
\begin{gathered}
  \mathsf{Spec}(x)\\[-0.2em]
  \Updownarrow\;\mathsf{Rep}(x,a)\\[-0.2em]
  \mathsf{ArrayWordState}(a,m)\\[-0.2em]
  \Updownarrow\;\mathsf{Memory}(m)\\[-0.2em]
  \mathsf{ConcreteBytesAndPointers}.
\end{gathered}
\]

여기서 \(\mathsf{Rep}\)는 결과 equality의 다른 이름이 아니라, 입력 배열이 어떤
수학 object를 나타내는지 정의하는 독립 predicate이다. 필요한 bridge lemma는
다음 범주로 나눈다.

\begin{itemize}
  \item NTT/CRT 배열과 \(\Rq\)의 coefficientwise polynomial 연산,
  \item W64/W32 값과 정수 사이의 범위 및 no-wraparound 관계,
  \item signed/unsigned 변환, arithmetic shift와 division by two,
  \item \(\HighBits\), \(\LowBits\), \(\LSB\), centered reduction과 rounding,
  \item byte offset, slice equality, encode/decode와 coefficient 순서,
  \item pointer separation, 변경 구간과 보존해야 할 prefix·suffix frame.
\end{itemize}

상수 범위나 memory-disjointness가 필요한 경우 theorem statement에 공개한다.
SMT가 우연히 concrete 상수를 풀었다는 사실을 representation theorem으로
대체하지 않는다.

\subsection{국소 Hoare 정리의 표준 형태}

각 actual procedure 정리는 다음 형태를 따른다.

\[
\begin{aligned}
\{\;&\mathsf{ArgBinding}
    \land \mathsf{Mode2Constants}
    \land \mathsf{InputRep}
    \land \mathsf{MemoryContract}\;\}\\
&\text{\proc{Actual.M._procedure}}\\
\{\;&\mathsf{ExactResult}
    \land \mathsf{OutputRep}
    \land \mathsf{BranchState}
    \land \mathsf{Frame}\;\}.
\end{aligned}
\]

\(\mathsf{ArgBinding}\)은 harness 인자와 procedure formal argument의 동일성을,
\(\mathsf{Mode2Constants}\)는 실제 \(q,n,k,\ell,\tau\) 및 길이를,
\(\mathsf{MemoryContract}\)는 유효 범위와 alias 조건을 고정한다. postcondition은
수학 결과뿐 아니라 반환 code, 실행된 branch, offset, 변경하지 않은 배열 구간을
함께 기술한다.

다음 항목은 원하는 결론을 가정하는 것이므로 원칙적으로 precondition에 두지
않는다.

\begin{itemize}
  \item encoder 또는 decoder가 성공했다는 가정,
  \item decoded object가 입력과 같다는 가정,
  \item copied trace나 decoder byte read가 encoder 출력과 같다는 가정,
  \item Sign 또는 Verify의 최종 \(reject\) 값,
  \item KeyGen의 최종 키 방정식이나 packing equality.
\end{itemize}

이 값들은 실제 하위 procedure의 postcondition과 bridge lemma에서 도출해야 한다.
도출할 수 없는 경우에는 theorem 이름과 claim 상태에 그 한계를 표시한다.

\subsection{대상별 증명 전략}

\begingroup
\small
\begin{longtable}{L{27mm}L{45mm}L{70mm}}
\toprule
대상 & 실제 증명 경계 & 핵심 증명 방법 \\
\midrule
\endhead
KeyGen &
\proc{_kp_m23_matrix}, \proc{_keypair_finalize_m23}의 focused actual sequence &
matrix/NTT 결과를 \textup{(KG-1)}로 refine하고 finalization에서
\textup{(KG-2)}를 얻는다. 배열로 구성한 \(A,s\)에 환 계산을 적용하여
\(As=qj\pmod{2q}\)를 도출하고 return/frame 상태를 보존한다. accepted parent의
guard, 길이와 packing lift는 7일 MINCORE 완료 조건에서 제외한다. \\
\addlinespace
Sign &
\proc{_sf_round_challenge_mode2}, \proc{_sf_z_check},
\proc{_sf_hint_mode2}의 순차 실행 &
raw \(\mu\) memory transcript와 challenge hash 입력을 연결하고, NTT challenge
product를 \(cs\)로 refine한다. fixed-point accumulator와 정수 squared norm의
범위를 증명하여 두 rejection comparison을 동치화하고, high-bit bridge로 hint
식을 얻는다. \\
\addlinespace
Verify &
canonical object 이후 \proc{_verify_prepare_z1_wprime}부터
\proc{_sign_verify_tail_m23}까지 &
byte/high-low decomposition에서 \(\widetilde z_1,w'\)를 복원하고, CRT matrix
결과를 \(u\)로 refine한다. 분자의 parity와 word 범위를 이용해 arithmetic shift를
정수 나눗셈으로 연결한 뒤, 반환 \(reject=0\)과 norm·challenge conjunction의
필요충분 동치를 증명한다. \\
\addlinespace
HBZ/rANS codec (기존 증거) &
actual encoder, suffix copy, actual decoder와 full wrapper; Week~16 확장 없음 &
encoder의 \texttt{segment\_matches}, copy의 \texttt{slice\_eq}, decoder의
pointwise read 조건을 adapter lemma로 연결한다. failure/success postcondition과
frame을 모두 유지한다. fixed all-6 결과는 종료한 실행의 실패만 배제하며
termination/losslessness를 제공하지 않는다고 표기한다. \\
\addlinespace
제한 합성 &
accepted KeyGen·Sign 결과와 actual Verify core &
동일 VK, prefix/context, message memory로 \(\mu=\widetilde\mu\)를 얻고 명시적
decoded-object identity로 signature object를 전달한다. KeyGen의 키 식과 Sign의
hint 식을 Verify 복원 lemma에 대입하여 최종 challenge equality와 norm 승인을
도출한다. \\
\bottomrule
\end{longtable}
\endgroup

\subsection{제어흐름, 종료성과 non-vacuity}

KeyGen과 Sign에는 rejection loop가 있으므로 일반 headline theorem은
\textbf{종료한 accepted execution}에 대한 Hoare partial correctness로 작성한다.
이는 해당 실행이 종료하면 postcondition이 성립한다는 뜻이지, 모든 입력에서
loop가 종료한다거나 기대 반복 횟수가 유한하다는 뜻이 아니다. loop 본문에서는
accepted branch에 필요한 invariant와 rejected iteration 뒤의 frame을 증명하고,
termination/losslessness는 별도 claim으로 관리한다.

codec처럼 실제 procedure가 실패할 수 있는 경우에는 결과를 다음처럼 분기한다.

\begin{itemize}
  \item failure branch: 반환 code, 미실행 후속 절차와 보존된 상태,
  \item success branch: 정확한 크기·trace·decoded equality와 frame.
\end{itemize}

success-conditioned partial correctness만으로는 성공 branch가 실제로 도달한다고
결론낼 수 없다. Week~15 결과는 all-zero HBZ의 \emph{종료한} 실행이 성공함을
보이지만 termination/losslessness를 보이지 않으므로 non-vacuous witness로
부르지 않는다. Verify core는 고정 길이 loop의 bounded 실행으로 다루며, 전체
parser나 malformed-input 경로를 포함하지 않으면 그 사실을 precondition에
명시한다.

\subsection{합성 정리와 완료 증거}

국소 정리는 caller가 실제로 제공하는 배열과 memory를 precondition에 대입할 수
있을 때만 합성한다. composition harness는 production 호출 순서를 보존하고,
procedure 사이에서 다음 상태를 명시적으로 전달한다.

\begin{itemize}
  \item KeyGen이 만든 VK/SK object와 \(As=qj\pmod{2q}\),
  \item Sign과 Verify가 공유하는 VK, prefix/context, message transcript,
  \item Sign의 \((z,h,c)\)와 동일 object를 나타내는 decoded \((x,v,h,c)\),
  \item 각 호출의 pointer binding, 변경 구간과 frame.
\end{itemize}

최종 목표는 다음 implication을 actual procedure composition의 postcondition으로
도출하는 것이다.

\[
\begin{gathered}
  \mathsf{TerminatedAcceptedKeyGen}
  \land \mathsf{TerminatedAcceptedSign}\\
  {}\land \mathsf{SameKeyAndMessageTranscript}
  \land \mathsf{DecodedSignatureObjectIdentity}\\
  {}\Longrightarrow \mathsf{ActualVerifyCore.reject}=0.
\end{gathered}
\]

정리는 EasyCrypt가 한 번 통과했다는 사실만으로 완료 처리하지 않는다. 각 신규
target의 fresh \proc{-no-eco} compile, 전체 aggregate compile, proof-hole·authored
axiom·debug declaration scan, extraction regeneration/hash 및 generated table drift,
read-only source root, safety 결과와 독립 audit를 함께 확인한다. theorem manifest,
claim ledger, 보고서와 논문이 같은 precondition과 claim status를 가질 때만
\cref{sec:artifact-plan}의 완료 기준을 만족한다.
